\section{Subarray Sums II}
\label{sec:cses-subarray-sums-ii}

\problemstatement{Given an array of $n$ integers (which may be negative)
and a target $x$, count the number of contiguous subarrays whose sum equals
$x$.}

\begin{patternbox}
With negatives, the window sum is no longer monotone, so two pointers fail.
Instead use \textbf{prefix sums plus a hash map}: a subarray $[l+1,r]$ has
sum $x$ iff $\textit{pre}[r]-\textit{pre}[l]=x$, so count, for each prefix,
how many earlier prefixes equal $\textit{pre}[r]-x$.
\end{patternbox}

\subsection*{Counting equal prefix differences}

Let $\textit{pre}[i]$ be the sum of the first $i$ elements. A subarray
ending at position $r$ with sum $x$ corresponds to an earlier prefix
$\textit{pre}[l]=\textit{pre}[r]-x$. Sweep left to right maintaining a hash
map from prefix-sum value to how many times it has occurred; at each step
add the map's count of $\textit{pre}[r]-x$ to the answer, then record
$\textit{pre}[r]$. Initialise the map with $\{0:1\}$ so subarrays starting
at index $0$ are counted.

\begin{ideabox}[Turn a Subarray Sum into a Prefix Difference]
Any subarray sum is a difference of two prefix sums, so ``sum equals $x$''
becomes ``two prefixes differ by $x$.'' A hash map of previously seen prefix
values counts those pairs in one pass, negatives and all.
\end{ideabox}

\subsection*{Algorithm}

\begin{enumerate}
  \item Map $\{0:1\}$, $\textit{pre}=0$, $\textit{count}=0$.
  \item For each element: $\textit{pre}\mathrel{+}=a[i]$; add
        $\textit{map}[\textit{pre}-x]$ to $\textit{count}$; increment
        $\textit{map}[\textit{pre}]$.
  \item Output $\textit{count}$.
\end{enumerate}

\subsection*{Dry run}

\dryrun{$a=[2,-1,3,1]$, $x=4$.}

\begin{itemize}
  \item Prefixes: $0,2,1,4,5$. Looking for earlier prefix
        $\textit{pre}-4$.
  \item At $\textit{pre}=4$: $\textit{pre}-4=0$ seen once $\Rightarrow$
        count $1$ ($[2,-1,3]$). At $\textit{pre}=5$: $5-4=1$ seen once
        $\Rightarrow$ count $2$ ($[3,1]$).
  \item Answer $2$.
\end{itemize}

\subsection*{C++ implementation}

\begin{lstlisting}
#include <bits/stdc++.h>
using namespace std;

int main() {
    int n; long long x;
    cin >> n >> x;
    unordered_map<long long,long long> seen;
    seen[0] = 1;                               // empty prefix

    long long pre = 0, count = 0;
    for (int i = 0; i < n; i++) {
        long long a; cin >> a;
        pre += a;
        auto it = seen.find(pre - x);          // earlier prefix that closes a sum-x subarray
        if (it != seen.end()) count += it->second;
        seen[pre]++;
    }
    cout << count << "\n";
    return 0;
}
\end{lstlisting}

\begin{complexitybox}
\textbf{Time Complexity.} $O(n)$ amortised with a hash map. \textbf{Space
Complexity.} $O(n)$ for the map.
\end{complexitybox}

\begin{takeawaybox}
Once negatives break monotonicity, prefix sums plus a hash map count
exact-sum subarrays in one pass -- the go-to when two pointers cannot apply.
This is the same machinery behind ``subarray with sum divisible by $n$''
(next section), swapping the key from a value to a residue.
\end{takeawaybox}
